\section{从账本看成化至正德的地方秩序}

本节按账本的时间序列观察地区、财政、边疆与社会问题。官修实录、诏令和奏疏强于保存中央选择记录的行动与日期；地方志、碑刻、契约、工程遗存及其他政权材料才可能校正具体地点的执行、损失与反应。因此，以下叙事不把行政分类、报送数字或动机判断直接等同于地方社会本身。

\subsection{流民、设治与山区社会}

\eventanchor{MM-1464-XIANZONG-ACCESSION}{明·1464（天顺八年）正月}
\paragraph{1464（天顺八年）正月：英宗去世，皇太子朱见深即位，是为宪宗} 账本首先限定我们能够确认的范围：死亡、即位与年号转换可靠；先朝遗诏和临终安排的细节须注意实录的礼制化记录。可供互读的材料包括《英宗实录》天顺八年正月条；《宪宗实录》成化元年即位条；《明史·宪宗本纪》。这些文本并非同一份现场证言，而是从朝廷、地方或跨政权的位置留下观察。事件之所以进入行政链，与；它带来的可追踪变化是成化朝得以在既有储位基础上开局，但宦官、后宫与边疆军务将构成新的权力场域。上述证据限度同样要求把日期、报数、执行与后果分开处理。

\eventanchor{ML-1465-001}{明·1465（成化元年）}
\paragraph{1465（成化元年）：蒙古诸部、朝鲜及西北绿洲的交涉在本年留下多方材料，应以具体使团和地名校正概称} 账本首先限定我们能够确认的范围：来源导向的年度桥接条目：本年材料可定位相关政务线索；具体月日、发文机关、地域和执行结果仍须逐项回查，不能以制度文本或奏报替代实际效果。可供互读的材料包括《宪宗实录》成化元年条；《明史·外国传》；《朝鲜王朝实录》、琉球《历代宝案》或海外档案。这些文本并非同一份现场证言，而是从朝廷、地方或跨政权的位置留下观察。事件之所以进入行政链，与；它带来的可追踪变化是它为后续册封、互市或海防安排留下文书与跨政权比较线索。桥接条目只标示可回查的年度问题与材料链；实录有中央选择性，崇祯期尤其需要奏疏、地方、朝鲜及满文材料交叉。

\eventanchor{ML-1465-002}{明·1465（成化元年）}
\paragraph{1465（成化元年）：地方灾异、赈济与水利条目为本年社会史提供入口，地方志的修纂层次需要说明} 账本首先限定我们能够确认的范围：来源导向的年度桥接条目：本年材料可定位相关政务线索；具体月日、发文机关、地域和执行结果仍须逐项回查，不能以制度文本或奏报替代实际效果。可供互读的材料包括《宪宗实录》成化元年条；《明史·五行志》；受灾府县地方志与奏疏。这些文本并非同一份现场证言，而是从朝廷、地方或跨政权的位置留下观察。事件之所以进入行政链，与；它带来的可追踪变化是赈济、迁徙和损失的实际范围仍须以受灾地区材料分别核验。桥接条目只标示可回查的年度问题与材料链；实录有中央选择性，崇祯期尤其需要奏疏、地方、朝鲜及满文材料交叉。

\eventanchor{MM-1465-DATENG-GORGE-CAMPAIGN}{明·1465（成化元年）}
\paragraph{1465（成化元年）：韩雍率军进攻大藤峡，镇压当地反抗并重整广西军政} 账本首先限定我们能够确认的范围：征讨、主帅和朝廷建置可由实录与传记互证；官文中的族名、战果和“平定”不等于当地社群已被完全控制。可供互读的材料包括《宪宗实录》成化元年广西条；《明史·韩雍传》；广西地方志与大藤峡相关碑刻。这些文本并非同一份现场证言，而是从朝廷、地方或跨政权的位置留下观察。事件之所以进入行政链，与；它带来的可追踪变化是明廷增设和调整广西防务、行政节点；山地社会的土地、迁徙和征役后果仍需地方材料追踪。上述证据限度同样要求把日期、报数、执行与后果分开处理。

\eventanchor{MM-1465-JINGXIANG-UPRISING}{明·1465（成化元年）}
\paragraph{1465（成化元年）：荆襄山地流民聚众反抗，朝廷调兵镇压并酝酿安置方案} 账本首先限定我们能够确认的范围：起事及朝廷处置可靠；“流民”“盗贼”是行政分类，参与者来源与规模不能由官军奏报直接推出。可供互读的材料包括《宪宗实录》成化元年荆襄条；《明史·原杰传》；郧阳、襄阳地方志。这些文本并非同一份现场证言，而是从朝廷、地方或跨政权的位置留下观察。事件之所以进入行政链，与；它带来的可追踪变化是军事镇压之外，朝廷转向以招抚、编户和设治处理人口流动；原杰抚治成为这一转折的关键。上述证据限度同样要求把日期、报数、执行与后果分开处理。

\eventanchor{ML-1466-001}{明·1466（成化二年）一月}
\paragraph{1466（成化二年）一月：明军击败并俘获侯大狗，但叛乱仍在继续} 账本首先限定我们能够确认的范围：编年候选的年序可由官修与后出材料交叉；具体月日、参与者、战果或数字必须回查所列材料，不能将单一叙事当作定论。可供互读的材料包括《宪宗实录》成化二年条；《明史·兵志》及相关列传；边镇、地方志或对方材料。这些文本并非同一份现场证言，而是从朝廷、地方或跨政权的位置留下观察。事件之所以进入行政链，与；它带来的可追踪变化是它改变了后续调兵、补给或地方秩序的条件；战果和伤亡须分开核对。译写候选以编年材料定位；实录记载反映朝廷选择，崇祯期须另以奏疏、地方及敌对政权材料交叉。

\subsection{边市、使团与东北}

\eventanchor{ML-1466-002}{明·1466（成化二年）}
\paragraph{1466（成化二年）：明军斩建州左卫东山} 账本首先限定我们能够确认的范围：编年候选的年序可由官修与后出材料交叉；具体月日、参与者、战果或数字必须回查所列材料，不能将单一叙事当作定论。可供互读的材料包括《宪宗实录》成化二年条；《明史·本纪》及相关列传；诏令、奏疏或会典版本。这些文本并非同一份现场证言，而是从朝廷、地方或跨政权的位置留下观察。事件之所以进入行政链，与；它带来的可追踪变化是它影响后续人事与政策议程；动机和派系标签不能由单一叙事裁定。译写候选以编年材料定位；实录记载反映朝廷选择，崇祯期须另以奏疏、地方及敌对政权材料交叉。

\eventanchor{ML-1466-003}{明·1466（成化二年）}
\paragraph{1466（成化二年）：湖南及川黔边境苗族叛乱遭镇压} 账本首先限定我们能够确认的范围：编年候选的年序可由官修与后出材料交叉；具体月日、参与者、战果或数字必须回查所列材料，不能将单一叙事当作定论。可供互读的材料包括《宪宗实录》成化二年条；《明史·本纪》及相关列传；诏令、奏疏或会典版本。这些文本并非同一份现场证言，而是从朝廷、地方或跨政权的位置留下观察。事件之所以进入行政链，与；它带来的可追踪变化是它影响后续人事与政策议程；动机和派系标签不能由单一叙事裁定。译写候选以编年材料定位；实录记载反映朝廷选择，崇祯期须另以奏疏、地方及敌对政权材料交叉。

\eventanchor{ML-1466-004}{明·1466（成化二年）}
\paragraph{1466（成化二年）：刘通在襄阳附近起兵，兵败} 账本首先限定我们能够确认的范围：编年候选的年序可由官修与后出材料交叉；具体月日、参与者、战果或数字必须回查所列材料，不能将单一叙事当作定论。可供互读的材料包括《宪宗实录》成化二年条；《明史·兵志》及相关列传；边镇、地方志或对方材料。这些文本并非同一份现场证言，而是从朝廷、地方或跨政权的位置留下观察。事件之所以进入行政链，与；它带来的可追踪变化是它改变了后续调兵、补给或地方秩序的条件；战果和伤亡须分开核对。译写候选以编年材料定位；实录记载反映朝廷选择，崇祯期须另以奏疏、地方及敌对政权材料交叉。

\eventanchor{MM-1466-YUAN-JIE-RESETTLEMENT}{明·1466（成化二年）}
\paragraph{1466（成化二年）：原杰奉命抚治荆襄，招抚、编置部分流民并核定屯种} 账本首先限定我们能够确认的范围：原杰奉命与安置措施可证；所报户口、田亩和安置成效是行政统计，不能等同全部流民的实际经历。可供互读的材料包括《宪宗实录》成化二年条；《明史·原杰传》；《郧阳府志》有关设治和流民记载。这些文本并非同一份现场证言，而是从朝廷、地方或跨政权的位置留下观察。事件之所以进入行政链，与；它带来的可追踪变化是招抚与编置使国家控制由临时征剿转向常设行政，为郧阳府设立和区域赋役重整铺路。上述证据限度同样要求把日期、报数、执行与后果分开处理。

\eventanchor{ML-1467-001}{明·1467（成化三年）}
\paragraph{1467（成化三年）：明朝鲜联军击败建州女真并杀死李曼珠} 账本首先限定我们能够确认的范围：编年候选的年序可由官修与后出材料交叉；具体月日、参与者、战果或数字必须回查所列材料，不能将单一叙事当作定论。可供互读的材料包括《宪宗实录》成化三年条；《明史·外国传》；《朝鲜王朝实录》、琉球《历代宝案》或海外档案。这些文本并非同一份现场证言，而是从朝廷、地方或跨政权的位置留下观察。事件之所以进入行政链，与；它带来的可追踪变化是它为后续册封、互市或海防安排留下文书与跨政权比较线索。译写候选以编年材料定位；实录记载反映朝廷选择，崇祯期须另以奏疏、地方及敌对政权材料交叉。

\eventanchor{ML-1467-002}{明·1467（成化三年）}
\paragraph{1467（成化三年）：科举、学校和法律条文在本年制度材料中可追踪，不能据规范文本推定州县实践一致} 账本首先限定我们能够确认的范围：来源导向的年度桥接条目：本年材料可定位相关政务线索；具体月日、发文机关、地域和执行结果仍须逐项回查，不能以制度文本或奏报替代实际效果。可供互读的材料包括《宪宗实录》成化三年条；《明史·选举志》或《艺文志》；文集、碑刻或书目材料。这些文本并非同一份现场证言，而是从朝廷、地方或跨政权的位置留下观察。事件之所以进入行政链，与；它带来的可追踪变化是它提供制度和传播史的锚点，不能据此推断社会整体接受。桥接条目只标示可回查的年度问题与材料链；实录有中央选择性，崇祯期尤其需要奏疏、地方、朝鲜及满文材料交叉。

\subsection{财政、交通与地方工程}

\eventanchor{ML-1467-003}{明·1467（成化三年）}
\paragraph{1467（成化三年）：土木危机后的边镇整顿、京师防务或复辟前后军政调度在本年材料中构成连续问题} 账本首先限定我们能够确认的范围：来源导向的年度桥接条目：本年材料可定位相关政务线索；具体月日、发文机关、地域和执行结果仍须逐项回查，不能以制度文本或奏报替代实际效果。可供互读的材料包括《宪宗实录》成化三年条；《明史·兵志》及相关列传；边镇、地方志或对方材料。这些文本并非同一份现场证言，而是从朝廷、地方或跨政权的位置留下观察。事件之所以进入行政链，与；它带来的可追踪变化是它改变了后续调兵、补给或地方秩序的条件；战果和伤亡须分开核对。桥接条目只标示可回查的年度问题与材料链；实录有中央选择性，崇祯期尤其需要奏疏、地方、朝鲜及满文材料交叉。

\eventanchor{MM-1467-JIANZHOU-CAMPAIGN}{明·1467（成化三年）}
\paragraph{1467（成化三年）：明军会同朝鲜军发动对建州女真据点的征讨，后世常称“成化犁庭”} 账本首先限定我们能够确认的范围：联合征讨发生可靠；“犁庭”、俘斩数字及其是否造成全面毁灭，均带有明朝与朝鲜军功叙事的夸张。可供互读的材料包括《宪宗实录》成化三年建州条；《朝鲜王朝实录》世祖三年北征条；《明史·女真传》。这些文本并非同一份现场证言，而是从朝廷、地方或跨政权的位置留下观察。事件之所以进入行政链，与；它带来的可追踪变化是部分聚落和首领网络遭受重创，但建州诸部并未消失；辽东边政仍依赖招抚、贸易与军事威慑并用。上述证据限度同样要求把日期、报数、执行与后果分开处理。

\eventanchor{ML-1468-001}{明·1468（成化四年）五月}
\paragraph{1468（成化四年）五月：蒙古人在固原起义} 账本首先限定我们能够确认的范围：编年候选的年序可由官修与后出材料交叉；具体月日、参与者、战果或数字必须回查所列材料，不能将单一叙事当作定论。可供互读的材料包括《宪宗实录》成化四年条；《明史·本纪》及相关列传；诏令、奏疏或会典版本。这些文本并非同一份现场证言，而是从朝廷、地方或跨政权的位置留下观察。事件之所以进入行政链，与；它带来的可追踪变化是它影响后续人事与政策议程；动机和派系标签不能由单一叙事裁定。译写候选以编年材料定位；实录记载反映朝廷选择，崇祯期须另以奏疏、地方及敌对政权材料交叉。

\eventanchor{ML-1468-002}{明·1468（成化四年）五月12日}
\paragraph{1468（成化四年）五月12日：越南千军占领广西凭祥边境小镇} 账本首先限定我们能够确认的范围：编年候选的年序可由官修与后出材料交叉；具体月日、参与者、战果或数字必须回查所列材料，不能将单一叙事当作定论。可供互读的材料包括《宪宗实录》成化四年条；《明史·外国传》；《朝鲜王朝实录》、琉球《历代宝案》或海外档案。这些文本并非同一份现场证言，而是从朝廷、地方或跨政权的位置留下观察。事件之所以进入行政链，与；它带来的可追踪变化是它为后续册封、互市或海防安排留下文书与跨政权比较线索。译写候选以编年材料定位；实录记载反映朝廷选择，崇祯期须另以奏疏、地方及敌对政权材料交叉。

\eventanchor{ML-1468-003}{明·1468（成化四年）}
\paragraph{1468（成化四年）：地方灾异、赈济与水利条目为本年社会史提供入口，地方志的修纂层次需要说明} 账本首先限定我们能够确认的范围：来源导向的年度桥接条目：本年材料可定位相关政务线索；具体月日、发文机关、地域和执行结果仍须逐项回查，不能以制度文本或奏报替代实际效果。可供互读的材料包括《宪宗实录》成化四年条；《明史·五行志》；受灾府县地方志与奏疏。这些文本并非同一份现场证言，而是从朝廷、地方或跨政权的位置留下观察。事件之所以进入行政链，与；它带来的可追踪变化是赈济、迁徙和损失的实际范围仍须以受灾地区材料分别核验。桥接条目只标示可回查的年度问题与材料链；实录有中央选择性，崇祯期尤其需要奏疏、地方、朝鲜及满文材料交叉。

\eventanchor{ML-1468-004}{明·1468（成化四年）}
\paragraph{1468（成化四年）：景泰、天顺时期的继承、人事与诏令链条可在本年材料中定位，政变责任不可由单一史书裁断} 账本首先限定我们能够确认的范围：来源导向的年度桥接条目：本年材料可定位相关政务线索；具体月日、发文机关、地域和执行结果仍须逐项回查，不能以制度文本或奏报替代实际效果。可供互读的材料包括《宪宗实录》成化四年条；《明史·本纪》及相关列传；诏令、奏疏或会典版本。这些文本并非同一份现场证言，而是从朝廷、地方或跨政权的位置留下观察。事件之所以进入行政链，与；它带来的可追踪变化是它影响后续人事与政策议程；动机和派系标签不能由单一叙事裁定。桥接条目只标示可回查的年度问题与材料链；实录有中央选择性，崇祯期尤其需要奏疏、地方、朝鲜及满文材料交叉。

\subsection{港口、军政与地方网络}

\eventanchor{ML-1469-001}{明·1469（成化五年）}
\paragraph{1469（成化五年）：固原蒙古叛乱被镇压} 账本首先限定我们能够确认的范围：编年候选的年序可由官修与后出材料交叉；具体月日、参与者、战果或数字必须回查所列材料，不能将单一叙事当作定论。可供互读的材料包括《宪宗实录》成化五年条；《明史·兵志》及相关列传；边镇、地方志或对方材料。这些文本并非同一份现场证言，而是从朝廷、地方或跨政权的位置留下观察。事件之所以进入行政链，与；它带来的可追踪变化是它改变了后续调兵、补给或地方秩序的条件；战果和伤亡须分开核对。译写候选以编年材料定位；实录记载反映朝廷选择，崇祯期须另以奏疏、地方及敌对政权材料交叉。

\eventanchor{ML-1469-002}{明·1469（成化五年）}
\paragraph{1469（成化五年）：景泰、天顺时期的继承、人事与诏令链条可在本年材料中定位，政变责任不可由单一史书裁断} 账本首先限定我们能够确认的范围：来源导向的年度桥接条目：本年材料可定位相关政务线索；具体月日、发文机关、地域和执行结果仍须逐项回查，不能以制度文本或奏报替代实际效果。可供互读的材料包括《宪宗实录》成化五年条；《明史·本纪》及相关列传；诏令、奏疏或会典版本。这些文本并非同一份现场证言，而是从朝廷、地方或跨政权的位置留下观察。事件之所以进入行政链，与；它带来的可追踪变化是它影响后续人事与政策议程；动机和派系标签不能由单一叙事裁定。桥接条目只标示可回查的年度问题与材料链；实录有中央选择性，崇祯期尤其需要奏疏、地方、朝鲜及满文材料交叉。

\eventanchor{ML-1469-003}{明·1469（成化五年）}
\paragraph{1469（成化五年）：土木危机后的边镇整顿、京师防务或复辟前后军政调度在本年材料中构成连续问题} 账本首先限定我们能够确认的范围：来源导向的年度桥接条目：本年材料可定位相关政务线索；具体月日、发文机关、地域和执行结果仍须逐项回查，不能以制度文本或奏报替代实际效果。可供互读的材料包括《宪宗实录》成化五年条；《明史·兵志》及相关列传；边镇、地方志或对方材料。这些文本并非同一份现场证言，而是从朝廷、地方或跨政权的位置留下观察。事件之所以进入行政链，与；它带来的可追踪变化是它改变了后续调兵、补给或地方秩序的条件；战果和伤亡须分开核对。桥接条目只标示可回查的年度问题与材料链；实录有中央选择性，崇祯期尤其需要奏疏、地方、朝鲜及满文材料交叉。

\eventanchor{ML-1469-004}{明·1469（成化五年）}
\paragraph{1469（成化五年）：军饷、漕运和户籍赋役的本年记录揭示恢复秩序的行政成本，账面额与实收须分列} 账本首先限定我们能够确认的范围：来源导向的年度桥接条目：本年材料可定位相关政务线索；具体月日、发文机关、地域和执行结果仍须逐项回查，不能以制度文本或奏报替代实际效果。可供互读的材料包括《宪宗实录》成化五年条；《明会典》相关条目；地方志、黄册/契约/河工材料。这些文本并非同一份现场证言，而是从朝廷、地方或跨政权的位置留下观察。事件之所以进入行政链，与；它带来的可追踪变化是其法定安排不等于地方执行，后续须以地域材料追踪负担与成效。桥接条目只标示可回查的年度问题与材料链；实录有中央选择性，崇祯期尤其需要奏疏、地方、朝鲜及满文材料交叉。

\eventanchor{ML-1470-001}{明·1470（成化六年）}
\paragraph{1470（成化六年）：辽东巡抚陈越攻击女真人，向女真使馆索要贿赂} 账本首先限定我们能够确认的范围：编年候选的年序可由官修与后出材料交叉；具体月日、参与者、战果或数字必须回查所列材料，不能将单一叙事当作定论。可供互读的材料包括《宪宗实录》成化六年条；《明史·本纪》及相关列传；诏令、奏疏或会典版本。这些文本并非同一份现场证言，而是从朝廷、地方或跨政权的位置留下观察。事件之所以进入行政链，与；它带来的可追踪变化是它影响后续人事与政策议程；动机和派系标签不能由单一叙事裁定。译写候选以编年材料定位；实录记载反映朝廷选择，崇祯期须另以奏疏、地方及敌对政权材料交叉。

\eventanchor{ML-1470-002}{明·1470（成化六年）}
\paragraph{1470（成化六年）：刘通叛军残部再次叛乱} 账本首先限定我们能够确认的范围：编年候选的年序可由官修与后出材料交叉；具体月日、参与者、战果或数字必须回查所列材料，不能将单一叙事当作定论。可供互读的材料包括《宪宗实录》成化六年条；《明史·本纪》及相关列传；诏令、奏疏或会典版本。这些文本并非同一份现场证言，而是从朝廷、地方或跨政权的位置留下观察。事件之所以进入行政链，与；它带来的可追踪变化是它影响后续人事与政策议程；动机和派系标签不能由单一叙事裁定。译写候选以编年材料定位；实录记载反映朝廷选择，崇祯期须另以奏疏、地方及敌对政权材料交叉。
